\chapter{الاحتمالات والمتغيرات العشوائية}\label{ch:g11:prob}

تنمذج الاحتمالات التجارب التي تكون نتيجتها غير مؤكدة: رمية نرد،
وسحب ورقة لعب، ولعبة تُلعب. ويرسي هذا الفصل المفردات على
المجموعات المنتهية، ثم يقدّم البطل المركزي في نظرية الاحتمالات كلها:
\emph{\omterm{def:g11:prob:rv}{المتغير العشوائي}} و\emph{أمله الرياضي}. وتُراجَع هذه المفاهيم
وتُعمَّق في \cref{ch:g12:randvar}، وتليها الاحتمالات
الشرطية في \cref{ch:g12:condprob}.

\section{الاحتمال على مجموعة منتهية}

\begin{definition}[النموذج الاحتمالي]\label{def:g11:prob:model}
تُنمذَج تجربة ذات عدد منته من الإمكانيات بواسطة \emph{مجموعة الإمكانيات}\index{مجموعة الإمكانيات} $\Omega = \{\omega_1, \dots, \omega_n\}$ (وهي
مجموعة النتائج الممكنة) مع \emph{\omterm{def:g10:proba:distribution}{احتمال}} $\P$ يرفق بكل
إمكانية $\omega_i$ عددًا $p_i \geq 0$، حيث $p_1 + \dots + p_n = 1$. و\emph{الحادثة}\index{حادثة} هي جزء $A \subseteq \Omega$، و
\[
\P(A) = \sum_{\omega_i \in A} p_i .
\]
وعندما تكون كل الإمكانيات متساوية \omterm{def:g10:proba:distribution}{الاحتمال} ($p_i = \frac1n$: وهو \emph{نموذج
التساوي في الاحتمال}\index{نموذج متساوي الاحتمال})،
\[
\P(A) = \frac{\text{عدد الإمكانيات في } A}
{\text{عدد الإمكانيات في } \Omega}.
\]
\end{definition}

\begin{example}\label{ex:g11:prob:twodice}
ارمِ حجري نرد يمكن تمييزهما: فإن $\Omega$ هي مجموعة الأزواج المرتبة
$(i, j)$ البالغ عددها $36$ حيث $1 \leq i, j \leq 6$، وكلها متساوية \omterm{def:g10:proba:distribution}{الاحتمال}. \omterm{def:g11:prob:model}{والحادثة}
$S$ = ``المجموع يساوي $7$'' تحوي $6$ أزواج، هي
$(1,6)$ و $(2,5)$ و $(3,4)$ و $(4,3)$ و $(5,2)$ و $(6,1)$، إذن
$\P(S) = \frac{6}{36} = \frac16$. أما \omterm{def:g11:prob:model}{الحادثة} ``المجموع يساوي $12$'' فتحوي
الزوج $(6,6)$ وحده: فاحتمالها $\frac{1}{36}$. واختيار \omterm{def:g11:prob:model}{مجموعة الإمكانيات}
\emph{المناسبة} --- أزواجًا مرتبة لا غير مرتبة --- هو ما يجعل
\omterm{def:g11:prob:model}{نموذج التساوي في الاحتمال} قابلًا للتطبيق.
\end{example}

\begin{proposition}[قواعد الاحتمال]\label{prop:g11:prob:rules}
من أجل حادثتين $A, B \subseteq \Omega$:
\[
\P(\emptyset) = 0, \qquad \P(\Omega) = 1, \qquad
\P(\bar A) = 1 - \P(A),
\]
\[
\P(A \cup B) = \P(A) + \P(B) - \P(A \cap B),
\]
حيث $\bar A$ هي \emph{\omterm{def:g10:proba:operations}{متممة}} $A$ (أي كل الإمكانيات التي ليست في $A$).
وخاصة، يكون $\P(A \cup B) = \P(A) + \P(B)$ عندما تكون $A$ و $B$
\emph{\omterm{def:g10:proba:operations}{متنافيتين}} ($A \cap B = \emptyset$).
\end{proposition}

\begin{proof}
لا تحوي \omterm{def:g11:prob:model}{الحادثة} الخالية أي إمكانية (فمجموع خالٍ)، أما $\Omega$ فتحويها
كلها (فالمجموع الكلي $1$). وكل إمكانية تنتمي إلى واحدة بالضبط من $A$
و $\bar A$، إذن $\P(A) + \P(\bar A) = 1$. وأما \omterm{def:g10:numbers:interunion}{الاتحاد}: فجمع المقادير $p_i$
على $A$ ثم على $B$ يعدّ إمكانيات $A \cap B$ مرتين؛
وطرح $\P(A \cap B)$ يصحّح هذا العدّ المزدوج.
\end{proof}

\begin{example}
في قسم، يدرس $60\%$ الموسيقى ($M$)، ويمارس $45\%$ رياضة ($S$)، ويجمع
$25\%$ بينهما. فنسبة من يفعل واحدًا على الأقل هي
$\P(M \cup S) = 0.60 + 0.45 - 0.25 = 0.80$، إذن $20\%$ لا يفعلون أيًا منهما
(\omterm{def:g10:proba:operations}{المتممة}).
\end{example}

\section{المتغيرات العشوائية}

\begin{definition}[المتغير العشوائي وتوزيعه]\label{def:g11:prob:rv}
\emph{المتغير العشوائي}\index{متغير عشوائي} على $\Omega$ هو \omterm{def:g11:func:function}{دالة}
$X$ ترفق بكل إمكانية عددًا. و\emph{توزيعه}\index{توزيع} هو
جدول قيمه الممكنة $x_1, \dots, x_k$ مع احتمالاتها
\[
\P(X = x_i) = \P\bigl(\{\omega : X(\omega) = x_i\}\bigr),
\qquad
\sum_{i=1}^k \P(X = x_i) = 1 .
\]
\end{definition}

\begin{example}\label{ex:g11:prob:game}
لعبة: ادفع $2$ أورو، وارمِ حجر نرد غير مغشوش، وخذ القيمة الظاهرة إن كانت
$5$ على الأقل، ولا شيء فيما عدا ذلك. وليكن $G$ هو \emph{الربح} (المقبوض
ناقص المدفوع). فالإمكانيات $1, 2, 3, 4$ تعطي $G = -2$؛ والإمكانية $5$
تعطي $G = 3$؛ والإمكانية $6$ تعطي $G = 4$. \omterm{def:g11:prob:rv}{وتوزيع} $G$:
\begin{center}
\begin{tabular}{c|ccc}
$g$ & $-2$ & $3$ & $4$\\
\hline\rule{0pt}{11pt}%
$\P(G = g)$ & $\dfrac46$ & $\dfrac16$ & $\dfrac16$
\end{tabular}
\end{center}
\end{example}

\begin{method}[بناء جدول التوزيع]\label{met:g11:prob:table}
اسرد القيم الممكنة للمتغير $X$؛ ومن أجل كل قيمة، اجمع الإمكانيات
التي تنتجها واجمع احتمالاتها؛ ثم تحقق من أن سطر
الاحتمالات مجموعه $1$.
\end{method}

\section{الأمل الرياضي}

\begin{definition}[الأمل الرياضي]\label{def:g11:prob:expectation}
\emph{الأمل الرياضي}\index{أمل رياضي} (أو \omterm{def:g10:stats:mean}{المتوسط}) للمتغير $X$ هو \omterm{def:g10:stats:mean}{متوسط}
قيمه مرجَّحة باحتمالاتها:
\[
\E(X) = \sum_{i=1}^{k} \P(X = x_i)\; x_i .
\]
\end{definition}

\begin{remark}
$\E(X)$ هو ما سيكون عليه \omterm{def:g10:stats:mean}{متوسط} عدد ضخم من التكرارات المستقلة
للتجربة --- والصياغة الدقيقة لهذه العبارة هي
قانون الأعداد الكبيرة، المبرهَن عليه في \cref{ch:g12:sums}. ولاحِظ التشابه
الشكلي مع \omterm{def:g10:stats:mean}{متوسط} \omterm{def:g10:stats:series}{سلسلة إحصائية} بالتواترات
(\cref{def:g11:stat:mean}): فالاحتمالات تلعب دور
التواترات النسبية $\frac{n_i}{n}$.
\end{remark}

\begin{example}\label{ex:g11:prob:gameexp}
من أجل لعبة \cref{ex:g11:prob:game}:
\[
\E(G) = \frac46 \times (-2) + \frac16 \times 3 + \frac16 \times 4
= \frac{-8 + 3 + 4}{6} = -\frac16 \approx -0.17 .
\]
فاللاعب يخسر وسطيًا نحو $17$ سنتًا في كل لعبة: فاللعبة
غير مواتية. وتُسمّى اللعبة \emph{عادلة} عندما يكون الربح المأمول $0$.
\end{example}

\begin{proposition}[الخطية]\label{prop:g11:prob:linearity}
من أجل كل عددين حقيقيين $a, b$:
\[
\E(aX + b) = a\,\E(X) + b .
\]
\end{proposition}

\begin{proof}
يأخذ المتغير $aX + b$ القيمة $a x_i + b$ \omterm{def:g10:proba:distribution}{باحتمال}
$\P(X = x_i)$، إذن
\[
\E(aX + b) = \sum_i \P(X = x_i)(a x_i + b)
= a \sum_i \P(X = x_i)\,x_i + b \underbrace{\sum_i \P(X =
x_i)}_{=\,1} = a\,\E(X) + b .
\]
\end{proof}

\section{التباين والانحراف المعياري}

\begin{definition}[التباين]\label{def:g11:prob:variance}
يقيس \emph{تباين}\index{تباين} المتغير $X$ تبدد قيمه
حول \omterm{def:g11:prob:expectation}{الأمل الرياضي} $m = \E(X)$:
\[
\V(X) = \E\bigl((X - m)^2\bigr)
= \sum_{i=1}^{k} \P(X = x_i)\,(x_i - m)^2,
\qquad
\sigma(X) = \sqrt{\V(X)} .
\]
\end{definition}

\begin{proposition}[الصيغة المختصرة]\label{prop:g11:prob:konig}
\[
\V(X) = \E(X^2) - \E(X)^2 .
\]
\end{proposition}

\begin{proof}
انشر $(x_i - m)^2 = x_i^2 - 2m\,x_i + m^2$ داخل المجموع:
\[
\V(X) = \sum_i \P(X = x_i)\,x_i^2
- 2m \sum_i \P(X = x_i)\,x_i + m^2\sum_i \P(X = x_i)
= \E(X^2) - 2m^2 + m^2 ,
\]
وهذا هو $\E(X^2) - m^2$. وهذه هي صيغة المتغيرات العشوائية
للمتطابقة المبرهَن عليها من أجل السلاسل في \cref{prop:g11:stat:shortcut}.
\end{proof}

\begin{proposition}[التحويل التآلفي]\label{prop:g11:prob:affine}
$\V(aX + b) = a^2\,\V(X)$ و $\sigma(aX + b) = \abs a\,\sigma(X)$.
\end{proposition}

\begin{proof}
حسب الخطية، ينحرف $aX + b$ عن أمله الرياضي $a m + b$ بمقدار
$(aX + b) - (am + b) = a(X - m)$: فكل انحراف يُضرب في $a$، ومنه
كل انحراف مربع يُضرب في $a^2$، وكذلك متوسطها المرجَّح. وقد
اختفى الانسحاب $b$ --- فانسحاب متغير لا يغيّر
تبدده.
\end{proof}

\begin{example}\label{ex:g11:prob:variance}
من أجل لعبة \cref{ex:g11:prob:game}:
$\E(G^2) = \frac46 \times 4 + \frac16 \times 9 + \frac16 \times 16
= \frac{16 + 9 + 16}{6} = \frac{41}{6}$، إذن
\[
\V(G) = \frac{41}{6} - \left(-\frac16\right)^2
= \frac{41}{6} - \frac{1}{36} = \frac{245}{36} \approx 6.8,
\qquad
\sigma(G) \approx 2.6 .
\]
فخسارة $0.17$ في كل لعبة وسطيًا تأتي مع تأرجحات حجمها النموذجي
$2.6$ أورو --- \omterm{def:g11:prob:expectation}{فالأمل الرياضي} بلا \omterm{def:g11:stat:variance}{انحراف معياري} لا يروي إلا نصف
الحكاية.
\end{example}

\begin{omfigure}
\begin{tikzpicture}
\begin{axis}[
  width=8.6cm, height=5.2cm,
  ybar, bar width=14pt,
  xlabel={$g$}, ylabel={$\P(G = g)$},
  xmin=-3.5, xmax=5.5, ymin=0, ymax=0.8,
  xtick={-2,3,4}, ytick={1/6, 4/6},
  yticklabels={$\tfrac16$,$\tfrac46$},
  tick label style={font=\small},
  label style={font=\small}]
  \addplot[fill=omDef!40, draw=omDef] coordinates {
    (-2,0.6667) (3,0.1667) (4,0.1667)};
  \draw[omThm, dashed, thick] (axis cs:-0.1667,0) --
    (axis cs:-0.1667,0.75);
  \node[omThm, font=\small, anchor=south west] at (axis cs:-0.05,0.62)
    {$\E(G)$};
\end{axis}
\end{tikzpicture}

{\small \omterm{def:g11:prob:rv}{توزيع} الربح $G$ في
\cref{ex:g11:prob:game}، وأمله الرياضي $\E(G) = -\frac16$ (بالمتقطع):
وهو نقطة توازن الكتل الاحتمالية.}
\end{omfigure}

\section{تمارين}

\begin{exercise}[$\star$]\label{exo:g11:prob:1}
تُسحب ورقة من رزمة عادية فيها $52$ ورقة. احسب \omterm{def:g10:proba:distribution}{احتمال}
سحب: قلبًا؛ ملكًا؛ قلبًا أو ملكًا؛ لا قلبًا ولا
ملكًا.
\end{exercise}

\begin{exercise}[$\star$]\label{exo:g11:prob:2}
يُرمى حجرا نرد غير مغشوشين ويُسجَّل مجموعهما $S$. وباستعمال
نموذج $36$ إمكانية في \cref{ex:g11:prob:twodice}، احسب $\P(S = 6)$
و $\P(S \geq 10)$ و $\P(S \text{ زوجي})$.
\end{exercise}

\begin{exercise}[$\star$]\label{exo:g11:prob:3}
\omterm{def:g11:prob:rv}{متغير عشوائي} \omterm{def:g11:prob:rv}{توزيعه}
\begin{center}
\begin{tabular}{c|cccc}
$x$ & $-1$ & $0$ & $2$ & $5$\\
\hline
$\P(X=x)$ & $0.3$ & $0.2$ & $0.4$ & $p$
\end{tabular}
\end{center}
أوجد $p$، ثم احسب $\E(X)$.
\end{exercise}

\begin{exercise}[$\star$]\label{exo:g11:prob:4}
يحقق $X$ العلاقتين $\E(X) = 3$ و $\V(X) = 4$. أعطِ \omterm{def:g11:prob:expectation}{الأمل الرياضي} \omterm{def:g11:prob:variance}{والتباين}
\omterm{def:g11:stat:variance}{والانحراف المعياري} للمتغير $Y = 2X - 5$.
\end{exercise}

\begin{exercise}[$\star\star$]\label{exo:g11:prob:5}
عجلة حظ فيها $10$ قطاعات متساوية: $5$ منها تُظهر $0$ أورو، و $3$ تُظهر
$2$ أورو، و $2$ تُظهر $5$ أورو. وتكلفة اللعب $2$ أورو. أعطِ
\omterm{def:g11:prob:rv}{توزيع} الربح $G$ (المكسب ناقص الرهان)، واحسب $\E(G)$، ثم
قرّر هل اللعبة مواتية.
\end{exercise}

\begin{exercise}[$\star\star$]\label{exo:g11:prob:6}
كيس فيه $3$ كرات حمراء و $2$ زرقاء؛ وتُسحب كرتان
\emph{بدون} إرجاع. وليكن $X$ عدد الكرات الحمراء المسحوبة.
ابنِ \omterm{def:g11:prob:rv}{توزيع} $X$ (بسرد $20$ سحبة مرتبة، أو باستعمال
\omterm{met:g10:proba:tree}{شجرة})، ثم احسب $\E(X)$.
\end{exercise}

\begin{exercise}[$\star\star$]\label{exo:g11:prob:7}
من أجل عجلة \cref{exo:g11:prob:5}، احسب $\V(G)$ و $\sigma(G)$.
\end{exercise}

\begin{exercise}[$\star\star$]\label{exo:g11:prob:8}
تبيع شركة تأمين عقدًا بمبلغ $80$ أورو. وباحتمال
$0.05$ يطالب الزبون بمبلغ $1000$ أورو، وباحتمال $0.01$ يطالب الزبون
بمبلغ $5000$ أورو، ولا شيء فيما عدا ذلك. وليكن $B$ ربح الشركة
من عقد واحد. أعطِ \omterm{def:g11:prob:rv}{توزيع} $B$، واحسب $\E(B)$، ثم
فسّر.
\end{exercise}

\begin{exercise}[$\star\star$]\label{exo:g11:prob:9}
سؤال متعدد الخيارات له $4$ أجوبة، واحد منها صحيح. والجواب الصحيح
يُعطي $+1$؛ وللتثبيط عن التخمين، يُعطي الجواب الخاطئ $x < 0$.
ويجيب تلميذ عشوائيًا. فمن أجل أي قيمة للعدد $x$ يكون
أمل نتيجة التلميذ معدومًا؟
\end{exercise}

\begin{exercise}[$\star\star$]\label{exo:g11:prob:10}
تُرمى قطعتا نقد غير مغشوشتين. وتقترح آن: ``إذا تطابقت القطعتان، أدفع لك
$1$ أورو؛ وإذا اختلفتا، تدفع لي $1$ أورو.'' احسب الربح المأمول
لكل لاعب. فهل اللعبة عادلة؟ والسؤال نفسه بثلاث قطع،
حيث تدفع آن $2$ أورو إذا تطابقت الثلاث.
\end{exercise}

\begin{exercise}[$\star\star\star$]\label{exo:g11:prob:11}
يانصيب يبيع $1000$ تذكرة بسعر $5$ أورو للواحدة. وتربح تذكرة واحدة
$2000$ أورو، وتربح خمس تذاكر $200$ أورو، وتربح عشرون تذكرة $25$
أورو.
\begin{enumerate}
  \item أعطِ \omterm{def:g11:prob:rv}{توزيع} مدخول المنظّم الإجمالي ناقص
        الجوائز، \omterm{def:g11:prob:rv}{وتوزيع} ربح لاعب واحد $G$.
  \item احسب $\E(G)$ وإجمالي الربح المأمول للمنظّم.
        وتحقق من اتساق الجوابين.
\end{enumerate}
\end{exercise}

\section{مسألة: اللعبة التي لم تُكمَل}

\begin{problem}[{مسألة نهاية الأسبوع --- باسكال وفيرما والقسمة
العادلة لرهان مقطوع: مولد الأمل الرياضي، وسحر
الخطية، وما لا يراه الأمل الرياضي}]
\label{pb:g11:prob:1}
في سنة 1654 طرح الفارس دو ميري --- وهو المقامر نفسه الذي فتحت
رهاناته على النرد المسألة الاحتمالية في الكتاب السابق ---
على بليز باسكال سؤالًا أعمق: لاعبان متساويا المهارة
يجب أن يوقفا سلسلتهما عند النتيجة $5$ إلى $3$، وستة
انتصارات تأخذ الرهان البالغ $64$ بستولًا. \emph{فكيف يُقسَم الرهان
قسمة عادلة؟} والرسائل المتبادلة بين باسكال وفيرما جوابًا عنه
أسّست المفهوم الذي هو قلب هذا الفصل: \omterm{def:g11:prob:expectation}{الأمل الرياضي}
(\cref{def:g11:prob:expectation}). وستعيد استخراج
جوابهما بالطريقتين، ثم تلقى حيلة الخطية البارعة وبقعة
\omterm{def:g11:prob:expectation}{الأمل الرياضي} العمياء الوحيدة.

\medskip
\noindent\textbf{الجزء الأول --- التمكّن من \omterm{def:g11:prob:expectation}{الأمل الرياضي}.}
\begin{enumerate}
  \item ارمِ حجر نرد غير مغشوش؛ وربحك هو (الوجه $- 3$) أورو. ابنِ
        جدول \omterm{def:g11:prob:rv}{توزيع} الربح $G$
        (\cref{met:g11:prob:table}) واحسب $\E(G)$.
  \item احسب $V(G)$ بالصيغة المختصرة
        (\cref{prop:g11:prob:konig}) ثم $\sigma(G)$ (برقمين
        عشريين).
  \item بلا أي جدول جديد، أعطِ $\E(2G - 1)$ و
        $V(2G - 1)$ (\cref{prop:g11:prob:affine}).
  \item كشك يتقاضى $m$ أورو مقابل رمي حجري نرد ويدفع
        لك المجموع بالأورو. فما السعر \emph{العادل}
        $m$؟ والكشك يتقاضى $8$: فما خسارتك المأمولة
        في اللعبة الواحدة؟
  \item حاسوب محمول قيمته $800$ أورو له \omterm{def:g10:proba:distribution}{احتمال} $5\,\%$ في
        التلف كل سنة. احسب قسط التأمين
        العادل، وربح المؤمِّن المأمول من العقد الواحد
        إذا تقاضى $60$ أورو.
\end{enumerate}

\medskip
\noindent\textbf{الجزء الثاني --- قسمة الرهان.}
النتيجة $5$--$3$، والأول الذي يبلغ $6$ يربح الرهان البالغ $64$ بستولًا، والجولات
رميات نقد عادلة.
\begin{enumerate}[resume]
  \item جوابان خاطئان مغريان: القسمة $5 : 3$
        (بالتناسب مع الجولات المربوحة)، أو إعطاء كل شيء
        للمتصدر. قل في جملة واحدة لكل منهما ما وجه
        الظلم فيه.
  \item طريقة فيرما: يحتاج اللاعب A إلى انتصار $1$ إضافي، ويحتاج B إلى
        $3$؛ فتخيّل أنهما يلعبان $3$ جولات إضافية بالضبط مهما كان.
        اسرد $8$ إمكانيات متساوية \omterm{def:g10:proba:distribution}{الاحتمال}، وحدّد من يربح
        السلسلة في كل منها، واستنتج القسمة العادلة للرهان البالغ $64$
        بستولًا.
  \item النقطة الدقيقة التي كان على فيرما الدفاع عنها: قد تتوقف
        اللعبة الحقيقية بعد جولة واحدة --- فلماذا يكون مشروعًا
        مع ذلك أن نستدل على $3$ جولات خيالية كاملة؟
  \item طريقة باسكال، بالرجوع من النهاية: احسب حصة A
        العادلة عند النتائج (الافتراضية) $5$--$5$، ثم
        $5$--$4$، ثم $5$--$3$، وفي كل مرة بأخذ \omterm{def:g10:stats:mean}{متوسط} المستقبلين
        المتساويي \omterm{def:g10:proba:distribution}{الاحتمال}. وقارن بالسؤال 7.
  \item السؤال نفسه عند النتيجة $4$--$3$: كم جولة خيالية،
        وأي الإمكانيات تعطي B السلسلة، وما هي
        القسمة العادلة؟
  \item صُغ المبدأ المشترك بين الطريقتين --- القيمة العادلة
        لمستقبل غير مؤكد هي \emph{أمله الرياضي}
        --- وسمِّ صناعتين حديثتين تقومان عليه بالكامل
        (السؤال 5 يبيّن إحداهما).
  \item تحقق من السلامة من أجل كل نتيجة: فسّر، بخطية
        \omterm{def:g11:prob:expectation}{الأمل الرياضي}
        (\cref{prop:g11:prob:linearity})، لماذا يجب أن يكون مجموع
        حصتي اللاعبين العادلتين دائمًا هو
        الرهان البالغ $64$ بستولًا بالضبط.
\end{enumerate}

\medskip
\noindent\textbf{الجزء الثالث --- حيلة الخطية البارعة.}
\begin{enumerate}[resume]
  \item استعد $\E(\text{مجموع حجري نرد}) = 7$ في سطر واحد
        بالخطية (\cref{prop:g11:prob:linearity})، بلا
        جدول من $36$ خانة. أما \emph{\omterm{def:g11:prob:variance}{تباين}} المجموع،
        فالقاعدة $V(X + Y) = V(X) + V(Y)$ تقتضي استقلال حجري
        النرد (وهو مقبول): فاحسبه.
  \item حفلة القبعات: يرمي $n$ ضيفًا قبعاتهم في
        كومة ويأخذ كل منهم واحدة عشوائيًا. وليكن $X$ عدد
        الضيوف الذين يستعيدون قبعتهم. فمن أجل ضيف واحد ثابت،
        ما \omterm{def:g10:proba:distribution}{احتمال} أن يستعيد قبعته
        هو؟ وبجمع هذه الفرص، وعددها $n$، (بالخطية!) احسب
        $\E(X)$. ويجب أن يفاجئك الجواب: فهو لا
        يتعلق بالعدد $n$.
  \item تحقق من $\E(X) = 1$ بالقوة الغاشمة من أجل $n = 2$ و
        $n = 3$ (اسرد توزيعات القبعات المتساوية \omterm{def:g10:proba:distribution}{الاحتمال}، وعددها $2$ ثم $6$،
        وخذ \omterm{def:g10:stats:mean}{متوسط} الأعداد).
  \item استعادات الضيوف بعيدة كل البعد عن الاستقلال (فإذا استعاد
        $n - 1$ ضيفًا قبعاتهم، فلا بد أن يستعيد الأخير
        قبعته أيضًا). فأين بالضبط \emph{لم} يقتضِ حساب السؤال 14
        الاستقلال؟ صُغ القوة الخارقة للخطية.
\end{enumerate}

\medskip
\noindent\textbf{الجزء الرابع --- ما لا يراه \omterm{def:g11:prob:expectation}{الأمل الرياضي}.}
\begin{enumerate}[resume]
  \item بطاقتا حظ ثمن كل منهما $5$ أورو: البطاقة A تربح $10$ أورو
        \omterm{def:g10:proba:distribution}{باحتمال} $\frac12$؛ والبطاقة B تربح $1\,000$ أورو
        \omterm{def:g10:proba:distribution}{باحتمال} $0.005$. تحقق من أن للربحين
        \omterm{def:g11:prob:expectation}{الأمل الرياضي} نفسه، ثم احسب التباينين. فأي
        بطاقة يشتريها الناس فعلًا، وماذا يشترون؟
  \item لعبة سان بطرسبورغ: ارمِ قطعة نقد حتى أول ظهر؛
        فإذا ظهر أول ظهر في الرمية $n$ تقبض
        $2^n$ أورو. احسب إسهام كل $n$ في
        \omterm{def:g11:prob:expectation}{الأمل الرياضي}، ثم \omterm{def:g11:prob:expectation}{الأمل الرياضي} نفسه. فهل تدفع
        $1\,000$ أورو للعب مرة واحدة؟ وماذا يعلّمنا هذا التصادم
        عن \omterm{def:g11:prob:expectation}{الأمل الرياضي}؟
  \item التأمين من جهة الزبون: أن تدفع $20$ أورو، أو
        أن تواجه خطرًا قدره $1\,\%$ في خسارة $1\,000$؟ قارن بين
        الأملين الرياضيين، وفسّر لماذا يبقى شراء التأمين
        عقلانيًا (فما الذي لا تستطيع أسرة واحدة أن
        ``تأخذ متوسطه'')؟
  \item الخاتمة --- سيرة \omterm{def:g11:prob:expectation}{الأمل الرياضي}، جملة واحدة لكل
        فصل منها: وُلد سنة 1654 ليقسم رهانًا؛ وقوته الخارقة،
        الخطية، التي تجمع الآمال الرياضية دون أن تسأل عن
        الاستقلال (القبعات)؛ وبقعته العمياء، الخطر، الذي يقيسه
        \omterm{def:g11:prob:variance}{التباين} (البطاقات، وسان بطرسبورغ، والتأمين)؛
        وتتويجه في السنة المقبلة، حين يفسّر قانون الأعداد
        الكبيرة بالضبط لماذا يربح بيت اللعب دائمًا.
\end{enumerate}
\end{problem}
